\section{Introduction}
\label{sec:introduction}
The superconformal index of a four-dimensional superconformal field theory was
introduced in \cite{Romelsberger2006,Kinney2007}. It is a generalisation of the
Witten index to superconformal field theories. Given mutually commuting charges
$G_i$ that commute with a chosen Poincaré supercharge $Q$ and its radial
quantisation adjoint $Q^\dagger$, the index is given by a trace over the Hilbert
space of the theory

\begin{equation}
  \mathcal I(\{\kappa_i\})
  =\operatorname{Tr}_{\mathcal H}
  \left[
    (-1)^F e^{-\beta\{Q,Q^\dagger\}}
    \prod_i \kappa_i^{G_i}
    \right],
  \label{eq:sc-index}
\end{equation}

where $F$ is the fermion number, $\beta$ an arbitrary constant, and the
$\kappa_i$ are fugacities for the remaining charges. The canonical reviews are
\cite{Rastelli:2016tbz,Gadde:2011uv,Gadde:2020yah}. Of particular importance is
the fact that the index tracks protected states, allowing exact results to be
obtained that are insensitive to coupling strength.

Here we consider 4d $\mathcal N=4$ super Yang-Mills theory. In the $SU(n)$ case
this theory appears on the conformal-field-theory side of the canonical duality
between type-IIB string theory on $AdS_5\times S^5$ and $\mathcal N=4$
SYM~\cite{Maldacena:1998}. In this case the superconformal index allows one to
count black hole microstates in the holographic dual.
More generally the index is useful as a quantitative test of field theory
dualities.

We are centrally interested in S-duality. This relates a theory with gauge
algebra $\mathfrak g$ to one with Langlands-dual algebra $\mathfrak g^{\vee}$.
For the $B_n$ and $C_n$ families,

\begin{equation}
  B_n^\vee=C_n,
  \qquad
  C_n^\vee=B_n,
\end{equation}

so the theories with gauge algebras $\mathfrak{so}(2n+1)$ and
$\mathfrak{sp}(2n)$ are exchanged. S-duality therefore predicts equality of
their superconformal indices. That equality is physically natural but is not
manifest in the corresponding matrix integrals.

Spiridonov and Vartanov expressed the superconformal indices of $\mathcal N=4$
SYM with simple gauge groups as elliptic hypergeometric
integrals~\cite{SpiridonovVartanov2012,Spiridonov:2020intro}. They clearly
articulated the S-duality property for these gauge groups as a conjectured
equality between two integrals. Moreover, they proved this identity in several
special cases but did not establish the generic finite-rank identity.

We will study a coarser version of the index called the Schur index. This is
constructed as a particular limit of the superconformal index which reduces the
number of fugacities. While containing less information, it becomes far more
tractable. For 4d theories with $\mathcal N=2$ superconformal symmetry, the
Schur index has several remarkable properties. For theories of class $\mathcal
S$, it can be computed using a $q$-deformed 2d Yang--Mills
theory~\cite{Gadde:2011ik}. The same quantity can also be identified with the
vacuum character of a corresponding chiral algebra
\cite{Beem:2013sza,Beem:2017ooy}. Of particular note, Bourdier, Drukker, and
Felix were able to derive a convergent closed expression for the index with
gauge group $U(N)$ \cite{Bourdier:2015wda}. A natural question is whether one
can derive a similar expression for the other classical groups, from which
S-duality should drop out.

There have been several attempts to make progress on this problem.

\begin{itemize}
  \item Pan and Peelaers developed a residue method for multivariate elliptic
        integrals and used it to obtain closed-form expressions for flavoured
        indices of several low-rank theories. In particular, they evaluated the
        $SO(5)$ theory which should be dual to the $Sp(4)$ theory
        \cite{Pan:2021ulr}.

  \item Hatsuda and Okazaki adapted the Fermi-gas approach of Bourdier,
        Drukker, and Felix. They obtained closed-form expressions for the Schur
        indices of 4d $\mathcal N=2^*$ SYM with unitary gauge groups at
        arbitrary rank \cite{Hatsuda:2022xdv}. Equivalently, from the $\mathcal
        N=4$ perspective, this restores the flavour fugacity for the adjoint
        hypermultiplet and gives the flavoured unitary-group analogue
        of~\cite{Bourdier:2015wda}.

  \item Guo, Li, Pan, and Wang extended the elliptic-integration method and
        derived a closed analytic expression for the fully flavoured Schur
        index of the $\mathcal N=4$ $SO(7)$ theory, the first group with
        non-trivial duality~\cite{Guo:2023qly}. They did not carry out the
        corresponding $Sp(6)$ calculation.

  \item Du, Huang, and Wang studied the unflavoured indices of the $B_n$,
        $C_n$, $D_n$, and $G_2$ theories using both character expansion and an
        extension of the Fermi-gas method \cite{Du:2023pbn}. For the cases they
        computed, they verified the equality of the $B_n$ and $C_n$ indices to
        high orders in the $q$-expansion. They also showed that the two
        families have the same stabilised large-$n$ index. However, they did
        not prove the $B_n/C_n$ equality at arbitrary rank.

  \item Ren and Huang recently studied a more general, three-parameter
        deformation of the Schur indices for the $B_n$, $C_n$, and $D_n$
        theories using Koornwinder and Macdonald polynomials
        \cite{RenHuang2026BCD}. For their coarser degeneration the index reduces
        to a known multivariable Askey--Wilson integral and the $B_n/C_n$
        S-duality becomes manifest.

\end{itemize}

In this paper, we will not seek closed-form expressions for these indices.
Instead we will demonstrate exact equality for the $B_n$ and $C_n$ indices at
the level of the integral representation. This paper is organised as follows:
Section~\ref{sec:schur} establishes the conventions and the two matrix
integrals. Section~\ref{sec:unflavoured} proves the unflavoured equality by
first expressing the indices in terms of a Fermi-gas partition function
representation, after which the two kernels are shown to be isospectral.
Section~\ref{sec:rains} restores the flavour fugacity and derives the equality
from a known elliptic hypergeometric integral transformation identity. We
conclude in Section~\ref{sec:conclusion}.

In Appendix~\ref{app:elliptic} we specify the theta- and
elliptic-gamma-function conventions. Appendices~\ref{app:mode-motivation} and
\ref{app:darboux} provide motivation and technical details for various aspects
of the unflavoured index calculation. Appendix~\ref{app:rains-limit} justifies
the boundary limit of the theorem used in the flavoured proof.

\textbf{Statement of AI use.} The construction reported on in this paper is
based on one originally generated by GPT-5.6 Sol and GPT-6 Astra at the author's
prompting. The author had periodically attempted to use the publicly available
models to solve the unflavoured version of this problem, typically to no avail.
Many of the ideas used in the final argument did emerge in responses from
earlier models but GPT-5.6 Sol was the first to put the pieces together and
solve the problem fully. GPT-6 Astra subsequently found the solution for the
flavoured case. The author refined the construction which is presented in this
paper. The manuscript was written by the author and we hope represents a clear
distillation of the arguments.
